\section{Konsep Kunci Publik}

Kriptografi kunci publik menggunakan sepasang kunci: \textbf{kunci publik} (dapat dibagikan) dan \textbf{kunci privat} (dijaga rahasia). Apa yang dienkripsi dengan kunci publik hanya dapat didekripsi dengan kunci privat, dan sebaliknya.

Keuntungan utama: tidak perlu saluran aman untuk pertukaran kunci. Siapa pun dapat mengenkripsi pesan untuk Bob menggunakan kunci publik Bob; hanya Bob yang dapat mendekripsi dengan kunci privatnya.

Kelemahan: lebih lambat dibanding kriptografi simetris. Oleh karena itu, sistem hibrida sering digunakan: kunci simetris dipertukarkan dengan kriptografi asimetris, lalu data dienkripsi dengan kunci simetris.
